% Written By Enoch Yu
% Downloaded from archive.enochyu.com

\documentclass{article}
\usepackage{graphicx}
\usepackage{amsmath,amsthm,amssymb}
\usepackage[font=small,labelfont=bf]{caption}
\usepackage{tikz}
\usepackage{pgfplots}
\pgfplotsset{compat=1.18}
\usetikzlibrary{calc, angles, quotes, shapes.geometric, decorations.pathreplacing}
\usepackage{tkz-euclide}
\usepackage[inline]{asymptote}
\usepackage{float}
\usepackage[margin=1in]{geometry}
\usepackage{gensymb}
\usepackage[normalem]{ulem}
\usepackage{hyperref}
\hypersetup{
    colorlinks=true,
    linkcolor=blue,
    filecolor=magenta,      
    urlcolor=cyan,
    pdftitle={Problem Set 32},
    pdfpagemode=FullScreen,
}
\usepackage{fancyhdr}
\pagestyle{fancy}
\fancyhead[R]{Enoch Yu}
\pagenumbering{gobble}
\usepackage{enumitem}
\newtheorem{theorem}{Theorem}[section]
\newtheorem{lemma}[theorem]{Lemma}
\newtheorem*{lemma*}{Lemma}
\newtheorem{sublemma}{Lemma}[section]
\newtheorem{proposition}{Proposition}
\newtheorem{corollary}{Corollary}[theorem]
\newtheorem{example}{Example}[section]
\newtheorem*{example*}{Example}
\newenvironment{solution}{\begin{trivlist}\item[]{\bf Solution}}{\qed \end{trivlist}}
\newcommand{\verteq}{\rotatebox{90}{$\;\;=\;\;$}}
\newcommand*\circled[1]{\tikz[baseline=(char.base)]{
            \node[shape=circle,draw,inner sep=1pt] (char) {#1};}}
\newcommand{\triangled}[1]{\tikz[baseline=(char.base)]{
            \node[shape=regular polygon, regular polygon sides=3, draw, inner sep=0.2pt] (char) {#1};}}

\title{Problem Set 32}
\author{Enoch Yu}
\date{July 2025}

\begin{document}

\section*{Problem}
Let $x = \sqrt{a^2 + a + 1} - \sqrt{a^2 - a + 1}$, where $a$ is real. Find all possible values of $x$.
\begin{solution}
Notice that because $a$ is a real number, discriminant could be used to show that $\sqrt{a^2 + a + 1} + \sqrt{a^2 - a + 1}$ is never zero. Thus, $x$ could be rewritten.
\[
x = \frac{2a}{\sqrt{a^2 + a + 1} + \sqrt{a^2 - a + 1}}
\]
Notice that when $a$ is greater than zero, the denominator is greater than $a$, leading to $x < 2$. When $a$ is zero, $x$ is simply zero. Furthermore, if $a$ is less than zero, then $x$ is merely a negative number. In other words, $x < 2$ for all values of $a$.
\begin{align*}
    x^2 + 2x\sqrt{a^2 - a + 1} + a^2 - a + 1 &= a^2 + a + 1 \\
    4x^2 \left( a^2 - a + 1 \right) &= 4a^2 - 4ax^2 + x^4 \\
    4x^2a^2 - 4a^2 &= x^4 - 4x^2 \\
    a^2 &= \frac{x^4 - 4x^2}{4x^2 - 4} = \frac{x^2(x^2 - 4)}{4(x^2 - 1)} \geq 0
\end{align*}
Notice that $\frac{x^2(x^2 - 4)}{4(x^2 - 1)}$ is positive \textit{iff} $x^2(x^2 - 4)(x^2 - 1)$ is positive. Meaning, $x^2$ is either in the interval $[0, 1]$ or $[4, \infty)$. However, the second interval is not applicable. Also, $\lim_{a\rightarrow\infty}\sqrt{a^2 + a + 1} - \sqrt{a^2 - a + 1} = 1$ and $\lim_{a\rightarrow-\infty}\sqrt{a^2 + a + 1} - \sqrt{a^2 - a + 1} = -1$. Therefore, $\boxed{-1 < x < 1}$.
\end{solution}

\section*{Problem}
Solve the equation $x^3 - 3x = \sqrt{x + 2}$ for all real values of $x$.
\begin{solution}
Notice that when $x = 2$, the equation satisfies. Therefore, intuitively, the cases for $x$ could be set.
\[
\begin{cases}
    x \leq 2 \implies \sqrt{x + 2} \geq x \\
    x > 2 \implies \sqrt{x + 2} < x
\end{cases}
\]
Observe the second equation.
\begin{align*}
    x^3 - 3x - \sqrt{x + 2} &= 0 \\
    x^3 - 3x - x &< 0 \\
    x(x^2 - 4) &< 0 \\
    \therefore x < -2 \text{ or } 0 &< x < 2
\end{align*}
Due to a contradiction, the fact that the roots of the equation are never greater than 2 could be deduced. Moreover, since $x$ is a real number, $-2 \leq x \leq 2$.
\\\\
\textcolor{red}{\textbf{GENERAL RULE OF THUMB: IF YOU SEE $-a \leq x \leq a$, IMMEDIATELY THINK OF TRIGONOMETRY!!!!}}
\\\\
Let $x = 2\cos\theta$. Therefore, the following equations could be deduced.
\begin{align*}
    8\cos^3\theta - 6\theta &= \sqrt{2\cos\theta + 2} \\
    \cos3\theta &= \sqrt{\frac{\cos\theta + 1}{2}} \\
    \cos3\theta &= \cos\frac{\theta}{2} \\
    \therefore x &= \boxed{2, 2\cos\frac{4\pi}{5}, 2\cos\frac{4\pi}{7}}
\end{align*}
\end{solution}

\end{document}
